% SPDX-License-Identifier: LicenseRef-MARKET-NC-Research-1.0
% GENERATED by: pandoc -s research/paper/manuscript.md -t latex (pandoc 3.6.3). Do not edit; edit manuscript.md. Not compiled locally (no TeX, D-003).
% Options for packages loaded elsewhere
\PassOptionsToPackage{unicode}{hyperref}
\PassOptionsToPackage{hyphens}{url}
\documentclass[
]{article}
\usepackage{xcolor}
\usepackage{amsmath,amssymb}
\setcounter{secnumdepth}{-\maxdimen} % remove section numbering
\usepackage{iftex}
\ifPDFTeX
  \usepackage[T1]{fontenc}
  \usepackage[utf8]{inputenc}
  \usepackage{textcomp} % provide euro and other symbols
\else % if luatex or xetex
  \usepackage{unicode-math} % this also loads fontspec
  \defaultfontfeatures{Scale=MatchLowercase}
  \defaultfontfeatures[\rmfamily]{Ligatures=TeX,Scale=1}
\fi
\usepackage{lmodern}
\ifPDFTeX\else
  % xetex/luatex font selection
\fi
% Use upquote if available, for straight quotes in verbatim environments
\IfFileExists{upquote.sty}{\usepackage{upquote}}{}
\IfFileExists{microtype.sty}{% use microtype if available
  \usepackage[]{microtype}
  \UseMicrotypeSet[protrusion]{basicmath} % disable protrusion for tt fonts
}{}
\makeatletter
\@ifundefined{KOMAClassName}{% if non-KOMA class
  \IfFileExists{parskip.sty}{%
    \usepackage{parskip}
  }{% else
    \setlength{\parindent}{0pt}
    \setlength{\parskip}{6pt plus 2pt minus 1pt}}
}{% if KOMA class
  \KOMAoptions{parskip=half}}
\makeatother
\usepackage{longtable,booktabs,array}
\usepackage{calc} % for calculating minipage widths
% Correct order of tables after \paragraph or \subparagraph
\usepackage{etoolbox}
\makeatletter
\patchcmd\longtable{\par}{\if@noskipsec\mbox{}\fi\par}{}{}
\makeatother
% Allow footnotes in longtable head/foot
\IfFileExists{footnotehyper.sty}{\usepackage{footnotehyper}}{\usepackage{footnote}}
\makesavenoteenv{longtable}
\usepackage{graphicx}
\makeatletter
\newsavebox\pandoc@box
\newcommand*\pandocbounded[1]{% scales image to fit in text height/width
  \sbox\pandoc@box{#1}%
  \Gscale@div\@tempa{\textheight}{\dimexpr\ht\pandoc@box+\dp\pandoc@box\relax}%
  \Gscale@div\@tempb{\linewidth}{\wd\pandoc@box}%
  \ifdim\@tempb\p@<\@tempa\p@\let\@tempa\@tempb\fi% select the smaller of both
  \ifdim\@tempa\p@<\p@\scalebox{\@tempa}{\usebox\pandoc@box}%
  \else\usebox{\pandoc@box}%
  \fi%
}
% Set default figure placement to htbp
\def\fps@figure{htbp}
\makeatother
\usepackage{svg}
\setlength{\emergencystretch}{3em} % prevent overfull lines
\providecommand{\tightlist}{%
  \setlength{\itemsep}{0pt}\setlength{\parskip}{0pt}}
\usepackage{bookmark}
\IfFileExists{xurl.sty}{\usepackage{xurl}}{} % add URL line breaks if available
\urlstyle{same}
\hypersetup{
  pdftitle={Trend followers, aggressive order flow and volatility clustering in a small simulated limit-order book: a preregistered paired-seed study},
  pdfauthor={Mohamed Bechir Kefi},
  hidelinks,
  pdfcreator={LaTeX via pandoc}}

\title{Trend followers, aggressive order flow and volatility clustering
in a small simulated limit-order book: a preregistered paired-seed
study}
\author{Mohamed Bechir Kefi}
\date{Draft, 2026-11-25}

\begin{document}
\maketitle

\textbf{Status.} Independent research manuscript / preprint draft, not
peer-reviewed. Not submitted anywhere. Author and affiliation metadata
are as approved by the owner (authorship decision C3); no institutional
affiliation, funding or endorsement is claimed.

\section{Abstract}\label{abstract}

Random-order (``zero-intelligence'') order flow in a continuous double
auction already reproduces much of the spread and diffusion behaviour of
real order books, and models with chartist traders have attributed
volatility clustering to trend-following feedback. We ask, in a small,
fully tested limit-order-book simulation with price-time priority and
cash and inventory constraints, whether replacing 30\% of random-order
agents by trend followers raises volatility clustering, measured by the
mean autocorrelation of absolute mid-price returns over lags 1 to 10
(VCI). The protocol, including an aggressiveness-matched placebo arm and
sensitivity analyses, was frozen before any confirmatory run; 480 runs
on 30 paired seeds were executed with no failure. The frozen verdict is
that the primary interpretation is invalidated: the frozen robustness
rule found a sign reversal, because the trend-follower contrast was
positive, with intervals excluding 0, at sampling intervals of 10 and
200 steps, while at the primary interval of 50 steps the paired VCI
difference was -0.007825 (95\% bootstrap interval {[}-0.0268,
0.01785{]}), inconclusive on its own. The rule takes its sign reference
from that near-zero primary estimate. As a secondary reading only,
taking the sign of H1 as the reference would make the same rule return
``qualified'' rather than ``invalidated''; both readings agree that the
answer depends on the sampling interval. At 50 steps, price-independent
aggressive agents raised VCI (placebo minus baseline 0.01751, {[}0.0135,
0.02148{]}), and trend followers added nothing beyond them (A1 minus
placebo -0.02534, {[}-0.04414, -0.0002837{]}). Secondary results: more
limit-order supply narrowed the spread, and trend followers deepened the
drawdown after a sell shock (both supported). Exploratory checks on
separate seeds indicate why a single interval can mislead here:
short-horizon VCI is shaped by bid-ask bounce and long-horizon VCI by
trend runs. The findings hold only within this model and make no
statement about real markets.

\section{Keywords}\label{keywords}

agent-based model; limit-order book; volatility clustering; trend
following; zero-intelligence traders; preregistration

\section{1. Introduction}\label{introduction}

Agent-based models of markets are often judged by whether they reproduce
``stylized facts'' of real returns, such as heavy tails and volatility
clustering {[}@Cont2001; @RatliffCrain2023{]}. Two traditions are
relevant here. In the first, random-order models of the continuous
double auction, with no strategic behaviour at all, account for much of
the cross-sectional variation of spreads and price diffusion in real
order books {[}@Farmer2005; @Smith2003{]}. In the second, models with
heterogeneous strategies attribute fat tails and clustering to the
interaction of trader types; in an order-driven model with
fundamentalist, chartist and noise components, the chartist component
was reported as mainly responsible for fat tails and clustering
{[}@Chiarella2009{]}.

A causal reading of the second tradition faces a confound. A trend
follower differs from a random-order trader in two ways at once: it
conditions its direction on past prices (feedback), and it sends a
different mix of market and limit orders (aggressiveness). Replacing
random-order traders by trend followers also removes liquidity
providers. A difference in clustering between the two markets therefore
does not, by itself, identify the feedback mechanism.

\textbf{Research question.} Within a simulated continuous double auction
with price-time priority, cash and inventory constraints and a fixed
order-arrival process, does replacing a share of random-order traders by
trend followers change the volatility clustering of mid-price returns,
and is any change attributable to price feedback rather than to
aggressiveness, the sampling interval or the tick size? Secondary
questions concern how liquidity supply and an exogenous shock change
spread, volume, volatility and drawdowns, and where the simulated
statistics sit relative to one real limit-order market.

\textbf{Hypotheses.} The single confirmatory hypothesis (H1, EXP-001) is
directional: the mean paired VCI difference between a market with 30\%
trend followers and an all-random-order market is positive. The
interpretation of H1 is guarded by preregistered sensitivity rules
(EXP-004).

\textbf{Contributions.} This is a replication-style study at small
scale, not a novelty claim. Its actual contributions are:

\begin{enumerate}
\def\labelenumi{\arabic{enumi}.}
\tightlist
\item
  a small, tested and deterministic order-book simulator (163 automated
  tests, including conservation, no negative balances, partial fills and
  price-time priority) with a reproducible runner and result schema;
\item
  a preregistered, paired-seed test of the trend-follower effect on
  volatility clustering against the random-order baseline, with an
  aggressiveness-matched placebo arm and sampling-interval and tick-size
  sensitivity rules frozen before any confirmatory output;
\item
  a complete report of the outcome, which is inconclusive at the primary
  interval and not robust under the frozen rule, together with the
  secondary results and a descriptive placement against one real
  order-book dataset (FI-2010);
\item
  an exploratory, post hoc sweep over sampling intervals on separate
  development seeds, showing that in this model VCI is entangled with
  bid-ask bounce at short horizons and with trend runs at long ones. It
  raises a question for a new preregistered study; it is not a finding
  about clustering.
\end{enumerate}

\section{2. Related work}\label{related-work}

This section synthesises only sources that were opened; for every source
except the FI-2010 paper and the full text of @RatliffCrain2023, only
the abstract or bibliographic record was read
(\texttt{research/literature\_review.md}). It is not a systematic
review.

\textbf{Random-order baselines.} @GodeSunder1993 showed that
budget-constrained random bids and offers reach near-full allocative
efficiency in a double auction, which established the zero-intelligence
trader as a baseline. @Smith2003 developed a statistical theory of the
continuous double auction under independent random order flow, with
predictions for volatility, spread and price impact. @Farmer2005 tested
a random order-placement model on London Stock Exchange data for eleven
stocks and report that, with one free parameter, it explains 96\% of the
variance of the bid-ask spread and 76\% of the variance of the price
diffusion rate across stocks. These results motivate the
all-random-order market (A0) as the null model of this study.

\textbf{Heterogeneous strategies and clustering.} @Chiarella2009 study
an order-driven market with fundamentalist, chartist and noise
components and report that the chartist strategy is mainly responsible
for fat tails and clustering in their artificial prices.
@LuxMarchesi1999 attribute scaling in a multi-agent (not order-book)
model to interactions among participants rather than to news.
@Chakraborti2011 review empirical facts and agent-based models,
including order-driven ones. @Byrd2019 describe ABIDES, a larger
simulator for market research; no code from it is used here.

\textbf{Stylized facts.} @Cont2001 catalogues stylized facts of asset
returns; its abstract refers to ``linear and nonlinear dependence of
returns in time'' without naming volatility clustering. The restatement
of Cont's list in @RatliffCrain2023, whose full text was read, includes
volatility clustering (``different measures of volatility display a
positive autocorrelation over several days'') and the slow decay of the
autocorrelation of absolute returns; those authors find the fact
supported in intraday Dow 30 data in both clock time and event time, and
caution that not all eleven facts hold in modern markets.

\textbf{Real order-book data.} FI-2010 {[}@Ntakaris2018{]} is a public
benchmark of NASDAQ Nordic limit-order-book data (five Helsinki stocks,
ten trading days in June 2010, one representation per 10 events, CC BY
4.0), built for mid-price forecasting. Here it is used only as a
descriptive comparison target, never for calibration.

\textbf{Relation to this study.} The study replicates, at small scale,
the qualitative question of @Chiarella2009 (do chartists produce
clustering?) and extends it with two controls that the opened abstracts
do not mention: a placebo arm that matches the trend followers' order
types and sizes but draws direction from a fair coin, and preregistered
sampling-interval and tick-size rules. Whether such placebo controls
exist in the wider literature was not searched systematically; no
priority is claimed.

\section{3. Background and formal
model}\label{background-and-formal-model}

The model is fixed in the frozen protocol
(\texttt{research/protocol.md}, sections 3 and 3.6).

\textbf{Market.} One asset on an integer price grid (tick 1 in the
primary configuration; initial price 1,000). A continuous double auction
with price-time priority: better price first, then earlier arrival.
Limit and market orders; market orders walk the book and may be
partially filled, with any remainder cancelled; marketable limit orders
execute up to their limit and rest the remainder. Each resting order
expires after \texttt{tau} = 2,000 steps. Trades occur at the resting
order's price. No fees, latency or hidden orders.

\textbf{Time.} Event time: at each step \texttt{t\ =\ 1..T} (\texttt{T}
= 100,000) one of 100 agents is drawn uniformly and acts once. The first
\texttt{B} = 10,000 steps are burn-in and are excluded from all metrics.

\textbf{Accounting.} Each agent holds cash
\texttt{c\_i\ \textgreater{}=\ 0} and inventory
\texttt{q\_i\ \textgreater{}=\ 0} (endowment 1,000,000 cash and 1,000
units); no short selling or borrowing. Orders are accepted only when
their worst-case cost is covered; totals of cash and inventory are
conserved exactly by every trade (integer arithmetic).

\textbf{Agents.} All agents share the order-size distribution (uniform
on 1..5 units) and the activation probability.

\begin{itemize}
\tightlist
\item
  \emph{Random-order (ZI):} with probability \texttt{p\_mkt} = 0.2 a
  market order on a random side; otherwise a limit order on a random
  side at an offset uniform on 0..10 ticks behind the best opposite
  quote.
\item
  \emph{Trend follower (TF):} signal
  \texttt{s\ =\ sign(m\_t\ -\ m\_\{t-L\})} with lookback \texttt{L} =
  500 steps on the mid-price \texttt{m}; no order when \texttt{s\ =\ 0};
  otherwise, in direction \texttt{s}, a market order with probability
  0.5, else a limit order at its own best quote.
\item
  \emph{Mean reverter (MR):} signal
  \texttt{s\ =\ -sign(m\_t\ -\ MA\_L(m)\_t)}, same order rule.
\item
  \emph{Placebo aggressive (PA):} the TF order-type and size
  distribution, direction from a fair coin independent of prices, acting
  on every activation.
\end{itemize}

\textbf{Randomness.} One \texttt{SeedSequence(seed)} is spawned into
named streams (activation, ZI decisions, strategic sizes, placebo coins,
shock, initial book). Arms that share a seed share the activation and ZI
streams (common random numbers), so arms are compared pairwise by seed.

\textbf{Primary metric.} With mid-price
\texttt{m\_t\ =\ (bid\_t\ +\ ask\_t)\ /\ 2} and returns sampled every
\texttt{dt} steps after burn-in,
\texttt{r\_k\ =\ ln\ m\_\{B\ +\ k\ dt\}\ -\ ln\ m\_\{B\ +\ (k-1)\ dt\}}
(1,800 returns per run at \texttt{dt} = 50), the volatility clustering
index is

\texttt{VCI\ =\ (1/10)\ sum\_\{l=1..10\}\ acf(\textbar{}r\textbar{},\ l)},

the mean sample autocorrelation of absolute returns at lags 1 to 10
(dimensionless; higher means more clustering). Secondary metrics
(protocol section 5) are realised volatility (SD of \texttt{r\_k},
log-return per \texttt{dt} steps), excess kurtosis, \texttt{acf(r,\ 1)},
mean quoted spread (ticks, over two-sided steps), depth at best (units),
volume per step (units), maximum drawdown of the mid (fraction),
one-sided fraction and rejected-order share.

\textbf{Complexity.} Price levels are held in heaps with lazy deletion
and each level is a first-in, first-out queue
(\texttt{src/market/book.py}), so a best-price query or an insertion is
amortised logarithmic in the number of levels and a market order costs
time proportional to the orders it fills; expiries use one more heap. A
run of 100,000 steps took on average between 1.4 and 2.05 seconds per
arm on a shared laptop (arm tables of the experiment reports); these are
not benchmark timings, because the machine was shared.

\section{4. Methods}\label{methods}

\subsection{4.1 Design and protocol
registration}\label{design-and-protocol-registration}

Five experiments were registered in
\texttt{research/experiment\_registry.csv}, each with its own protocol
(\texttt{research/experiments/EXP-00N/protocol.md}). The owner froze
protocol v1.0 on 2026-11-24 after two rounds of amendments
(\texttt{research/freeze.json}, decision log D-010 and D-011); the
freeze record holds the SHA-256 of every protocol and config, and the
runner refuses any config whose hash differs. The analysis module
(\texttt{market.frozen\_analysis}) was written, tested and committed
before any confirmatory output existed (D-013). Its implementation
choices where the frozen text was silent are listed in section 4.5 and
were reviewed in M4 (section 6.4).

\begin{longtable}[]{@{}
  >{\raggedright\arraybackslash}p{(\linewidth - 6\tabcolsep) * \real{0.2500}}
  >{\raggedright\arraybackslash}p{(\linewidth - 6\tabcolsep) * \real{0.2500}}
  >{\raggedright\arraybackslash}p{(\linewidth - 6\tabcolsep) * \real{0.2500}}
  >{\raggedright\arraybackslash}p{(\linewidth - 6\tabcolsep) * \real{0.2500}}@{}}
\toprule\noalign{}
\begin{minipage}[b]{\linewidth}\raggedright
EXP
\end{minipage} & \begin{minipage}[b]{\linewidth}\raggedright
Question
\end{minipage} & \begin{minipage}[b]{\linewidth}\raggedright
Arms
\end{minipage} & \begin{minipage}[b]{\linewidth}\raggedright
Runs
\end{minipage} \\
\midrule\noalign{}
\endhead
\bottomrule\noalign{}
\endlastfoot
EXP-001 & Do trend followers raise VCI against the random-order
baseline? (confirmatory) & A0 100 ZI; A1 70 ZI + 30 TF; A2 70 ZI + 30
MR; A3 70 ZI + 15 TF + 15 MR & 120 \\
EXP-002 & Does more ZI limit supply narrow the spread, and does
liquidity modulate the TF effect on VCI? & ZI market-order probability
0.35, 0.2, 0.1 crossed with ZI-only and ZI+TF & 180 \\
EXP-003 & After one market sell of 50\% of resting bid depth at step
50,000, do mixtures differ in drawdown? (optional) & S-ZI, S-ZI+TF,
S-ZI+MR & 90 \\
EXP-004 & Is any EXP-001 effect attributable to feedback rather than
aggressiveness, sampling interval or tick? & A4 70 ZI + 30 PA; A0 and A1
at tick 2; reanalysis of EXP-001 at \texttt{dt} 10 and 200 & 90 \\
EXP-005 & Where do the EXP-001 arms sit against FI-2010 on VCI, realised
volatility and relative spread? (descriptive) & EXP-001 runs against
reference ranges & 0 new \\
\end{longtable}

\subsection{4.2 Simulation model and
implementation}\label{simulation-model-and-implementation}

The engine (\texttt{src/market/}) implements the model of section 3 in
Python with integer prices and quantities. Invariants (conservation of
cash and inventory, no negative balance, price-time priority) are
checked at the end of every run, and a violation fails the run. Before
the confirmatory runs, 154 tests passed, all frozen hashes matched, and
the smoke matrix reproduced the earlier smoke run bit for bit (journal
J-0005); 13 of 13 metric checks on hand-computed cases passed
(\texttt{scripts/validate\_metrics.py}). At the confirmatory revision
the suite had 163 tests; the M4 revision adds three tests of the M4
scripts, for 166 in total, all passing (\texttt{progress.md},
2026-11-25).

Rules fixed during implementation and folded into the frozen protocol
(section 3.6) include: self-trades are allowed and accounting-neutral; a
market buy is capped by free cash; the initial book has 20 ZI limit
orders per side; prices have a floor of one tick.

\subsection{4.3 Baselines and controls}\label{baselines-and-controls}

The baseline is the all-random-order market A0, the simplest market in
which the same institutions operate. Three controls guard the
interpretation of EXP-001:

\begin{itemize}
\tightlist
\item
  \emph{Placebo aggressiveness (A4):} 30 PA agents replace 30 ZI agents.
  PA matches TF order types, sizes and limit placement, and acts on
  every activation; TF skips only when its signal is exactly 0, which
  happened in 0.4518\% of A1 activations on average (maximum 1.251\%),
  below the preregistered 5\% threshold for the match. A4 keeps the same
  number of ZI liquidity providers as A1, so A1 minus A4 isolates the
  direction rule from both aggressiveness and liquidity removal.
\item
  \emph{Sampling interval:} the EXP-001 contrast recomputed at
  \texttt{dt} 10 and 200 steps.
\item
  \emph{Tick size:} A0 and A1 on a grid of 2 price units.
\end{itemize}

\subsection{4.4 Seeds, resource budget and statistical
analysis}\label{seeds-resource-budget-and-statistical-analysis}

The experimental unit is one run (one arm, one seed). Confirmatory seeds
are 1001 to 1030 for every arm (30 per arm, chosen before the freeze for
80\% power at a standardised paired effect \texttt{d\_z} of about 0.53).
Development seeds 1 to 999 were used only for exploratory work. Every
contrast is the mean over seeds of per-seed paired differences
(treatment minus control), with a 95\% percentile bootstrap interval
(10,000 resamples, bootstrap seed 20261124) and \texttt{d\_z} = mean /
SD of the differences. A directional hypothesis is supported if the
interval lies above 0, contradicted if it lies below 0 and inconclusive
otherwise; an inconclusive result is never read as equivalence. Declared
secondary contrasts are reported with Holm-adjusted intervals within
each experiment. Runs flagged as degenerate (one-sided book in more than
5\% of steps, or fewer than 50 non-zero returns) are kept in the primary
analysis and removed, with their paired partners, in a sensitivity
analysis.

All computation ran on one laptop (Apple M1 Pro, 10 cores) at zero cost,
with two worker processes. The four batches took 111.3, 148.8, 73.3 and
68.8 seconds of wall time (\texttt{run\_manifest.json} per experiment).

\subsection{4.5 Exploratory versus confirmatory work and analysis
choices}\label{exploratory-versus-confirmatory-work-and-analysis-choices}

Confirmatory: the 480 runs at code revision \texttt{26266f2} and the
EXP-005 comparison on them. Exploratory: the smoke and pilot runs before
the freeze (J-0002, J-0003), the FI-2010 reference fit before the freeze
(J-0004), and the M4 checks of section 5.6 (journal J-0011), which used
development seeds 101 to 130 or read committed outputs only. No
exploratory result changes a frozen verdict.

Choices made by the AI assistant before the data where the frozen text
was silent (D-013): (1) Holm intervals use a two-sided bootstrap p-value
and a step-down level \texttt{1\ -\ 0.05\ /\ (m\ -\ j\ +\ 1)} for the
contrast of rank \texttt{j}; (2) the EXP-001 secondary family has 30
contrasts (A2 - A0, A3 - A0, A1 - A2 on ten metrics); (3) the EXP-002
family has 16 contrasts (liquidity levels within composition plus the
interaction); (4) the EXP-003 family has one contrast, and ``spread
widening'' was not measured because the protocol gives no formula; (5)
degenerate-run sensitivity drops the whole seed from a contrast; (6) in
the EXP-004 robustness rule, ``same sign'' means the sign of the EXP-001
primary point estimate.

\section{5. Results}\label{results}

All 480 confirmatory runs completed with status \texttt{ok}; none
failed, was dropped or was retried. Every run passed the end-of-run
invariant checks. Table 1 sets each frozen prediction beside its
observation.

\textbf{Table 1.} Expected versus observed for every registered
experiment. Source: \texttt{results/confirmatory/EXP-00N/analysis.json}
and \texttt{EXP-005/comparison.json}.

\begin{longtable}[]{@{}
  >{\raggedright\arraybackslash}p{(\linewidth - 6\tabcolsep) * \real{0.2500}}
  >{\raggedright\arraybackslash}p{(\linewidth - 6\tabcolsep) * \real{0.2500}}
  >{\raggedright\arraybackslash}p{(\linewidth - 6\tabcolsep) * \real{0.2500}}
  >{\raggedright\arraybackslash}p{(\linewidth - 6\tabcolsep) * \real{0.2500}}@{}}
\toprule\noalign{}
\begin{minipage}[b]{\linewidth}\raggedright
Experiment
\end{minipage} & \begin{minipage}[b]{\linewidth}\raggedright
Prediction (frozen)
\end{minipage} & \begin{minipage}[b]{\linewidth}\raggedright
Observation (95\% percentile bootstrap interval)
\end{minipage} & \begin{minipage}[b]{\linewidth}\raggedright
Verdict under the frozen rule
\end{minipage} \\
\midrule\noalign{}
\endhead
\bottomrule\noalign{}
\endlastfoot
EXP-001 (primary, confirmatory) & VCI(A1) - VCI(A0) \textgreater{} 0
(trend followers raise clustering) & -0.007825 {[}-0.0268, 0.01785{]},
d\_z -0.121, 30 pairs & inconclusive (H1 not supported) \\
EXP-002 (primary) & spread(low-ZI) - spread(high-ZI) \textgreater{} 0 &
1.632 price units {[}1.624, 1.641{]}, d\_z 69.3, 30 pairs & supported \\
EXP-002 (secondary, Holm) & interaction on VCI, nondirectional & 0.0555,
Holm interval {[}0.03614, 0.08209{]} & nonzero (secondary) \\
EXP-003 (primary, optional) & drawdown(S-ZI+TF) - drawdown(S-ZI)
\textgreater{} 0 & 0.02545 {[}0.01976, 0.03148{]}, d\_z 1.52, 30 pairs &
supported \\
EXP-004 placebo & VCI(A1) - VCI(A4) \textgreater{} 0 (feedback beyond
aggressiveness) & -0.02534 {[}-0.04414, -0.0002837{]} & contradicted;
attribution rule not triggered (EXP-001 inconclusive) \\
EXP-004 robustness rule & same sign and interval excluding 0 at dt 10,
dt 200 and tick 2 & vci\_dt10: 0.02481 {[}0.0153, 0.03842{]};
vci\_dt200: 0.0915 {[}0.04287, 0.1489{]}; tick 2: A1-tick2 minus
A0-tick2: -0.005427 {[}-0.03144, 0.02572{]} & invalidated (a sign
reversal with an interval excluding 0) \\
EXP-005 (a) & held-out day: VCI \textgreater{} 0 in each of 5 stocks &
-0.001241, 0.03691, 0.05438, 0.03336, 0.07877 & not supported (one stock
below 0) \\
EXP-005 (b) & no directional expectation (descriptive placement) & see
Table 7 & descriptive \\
\end{longtable}

\subsection{5.1 EXP-001: agent mixtures and volatility
clustering}\label{exp-001-agent-mixtures-and-volatility-clustering}

\textbf{Table 2.} EXP-001 arm summary: mean (SD) over 30 runs per arm.
Source: \texttt{results/confirmatory/EXP-001/tables/arm\_summary.csv}.

\begin{longtable}[]{@{}
  >{\raggedright\arraybackslash}p{(\linewidth - 16\tabcolsep) * \real{0.1111}}
  >{\raggedright\arraybackslash}p{(\linewidth - 16\tabcolsep) * \real{0.1111}}
  >{\raggedright\arraybackslash}p{(\linewidth - 16\tabcolsep) * \real{0.1111}}
  >{\raggedright\arraybackslash}p{(\linewidth - 16\tabcolsep) * \real{0.1111}}
  >{\raggedright\arraybackslash}p{(\linewidth - 16\tabcolsep) * \real{0.1111}}
  >{\raggedright\arraybackslash}p{(\linewidth - 16\tabcolsep) * \real{0.1111}}
  >{\raggedright\arraybackslash}p{(\linewidth - 16\tabcolsep) * \real{0.1111}}
  >{\raggedright\arraybackslash}p{(\linewidth - 16\tabcolsep) * \real{0.1111}}
  >{\raggedright\arraybackslash}p{(\linewidth - 16\tabcolsep) * \real{0.1111}}@{}}
\toprule\noalign{}
\begin{minipage}[b]{\linewidth}\raggedright
Arm
\end{minipage} & \begin{minipage}[b]{\linewidth}\raggedright
Runs
\end{minipage} & \begin{minipage}[b]{\linewidth}\raggedright
VCI
\end{minipage} & \begin{minipage}[b]{\linewidth}\raggedright
Realised volatility
\end{minipage} & \begin{minipage}[b]{\linewidth}\raggedright
Excess kurtosis
\end{minipage} & \begin{minipage}[b]{\linewidth}\raggedright
acf(r, 1)
\end{minipage} & \begin{minipage}[b]{\linewidth}\raggedright
Mean spread (ticks)
\end{minipage} & \begin{minipage}[b]{\linewidth}\raggedright
Depth at best (units)
\end{minipage} & \begin{minipage}[b]{\linewidth}\raggedright
Max drawdown (fraction)
\end{minipage} \\
\midrule\noalign{}
\endhead
\bottomrule\noalign{}
\endlastfoot
A0 & 30 & 0.02564 (0.008298) & 0.001302 (3.39e-05) & 0.1309 (0.1118) &
-0.3786 (0.02223) & 1.792 (0.01671) & 71.1 (1.992) & 0.01051
(0.002913) \\
A1 & 30 & 0.01781 (0.06483) & 0.0005241 (0.0001892) & 5.61 (20.02) &
-0.00584 (0.09718) & 1.3 (0.02458) & 109.4 (5.013) & 0.2004 (0.05875) \\
A2 & 30 & 0.0435 (0.01728) & 0.001391 (3.854e-05) & 0.03735 (0.1185) &
-0.4887 (0.01567) & 2.766 (0.02813) & 246 (21.96) & 0.0057
(0.0008209) \\
A3 & 30 & 0.05315 (0.01569) & 0.001151 (3.91e-05) & 0.7721 (0.2274) &
-0.3622 (0.02214) & 1.955 (0.02292) & 73.42 (2.922) & 0.01327
(0.003833) \\
\end{longtable}

The primary contrast VCI(A1) - VCI(A0) is -0.007825, 95\% interval
{[}-0.0268, 0.01785{]}, \texttt{d\_z} = -0.121, 30 pairs (Figure 1).
Under the frozen rule H1 is not supported: the result is inconclusive,
not evidence of no effect. No run in A0 or A1 was flagged, so the
degenerate-run sensitivity equals the primary analysis.

The main feature of Table 2 is dispersion. The SD of VCI across seeds is
0.06483 in A1 against 0.008298 in A0. Per-seed A1 values range from
-0.02913 to 0.3079, and 16 of 30 A1 runs have a negative VCI, while
every A0 run lies between 0.01082 and 0.0426 (exploratory summary E5 of
the committed runs). The study was powered for \texttt{d\_z} of about
0.53; the observed \texttt{\textbar{}d\_z\textbar{}} is about 0.12.

Trend followers change the market's levels strongly (Table 2): realised
volatility at \texttt{dt} = 50 falls from 0.001302 to 0.0005241, the
mean spread from 1.792 to 1.3 ticks, depth at best rises from 71.1 to
109.4 units, and the mean maximum drawdown rises from 0.01051 to 0.2004
of the running peak. Mean reverters (A2) and the mixture (A3) raise VCI
relative to A0, with Holm intervals excluding 0 (Table 3).

\textbf{Table 3.} Selected EXP-001 secondary contrasts (Holm family of
30; full table in the EXP-001 report).

\begin{longtable}[]{@{}
  >{\raggedright\arraybackslash}p{(\linewidth - 12\tabcolsep) * \real{0.1429}}
  >{\raggedright\arraybackslash}p{(\linewidth - 12\tabcolsep) * \real{0.1429}}
  >{\raggedright\arraybackslash}p{(\linewidth - 12\tabcolsep) * \real{0.1429}}
  >{\raggedright\arraybackslash}p{(\linewidth - 12\tabcolsep) * \real{0.1429}}
  >{\raggedright\arraybackslash}p{(\linewidth - 12\tabcolsep) * \real{0.1429}}
  >{\raggedright\arraybackslash}p{(\linewidth - 12\tabcolsep) * \real{0.1429}}
  >{\raggedright\arraybackslash}p{(\linewidth - 12\tabcolsep) * \real{0.1429}}@{}}
\toprule\noalign{}
\begin{minipage}[b]{\linewidth}\raggedright
Contrast
\end{minipage} & \begin{minipage}[b]{\linewidth}\raggedright
Metric
\end{minipage} & \begin{minipage}[b]{\linewidth}\raggedright
Mean diff.
\end{minipage} & \begin{minipage}[b]{\linewidth}\raggedright
d\_z
\end{minipage} & \begin{minipage}[b]{\linewidth}\raggedright
95\% interval
\end{minipage} & \begin{minipage}[b]{\linewidth}\raggedright
Holm interval (D-013)
\end{minipage} & \begin{minipage}[b]{\linewidth}\raggedright
Holm p
\end{minipage} \\
\midrule\noalign{}
\endhead
\bottomrule\noalign{}
\endlastfoot
A2 - A0 & vci & 0.01786 & 0.851 & {[}0.01049, 0.02521{]} & {[}0.005891,
0.02904{]} & \textless{} 0.0001 \\
A2 - A0 & realised\_volatility & 8.95e-05 & 1.89 & {[}7.261e-05,
0.0001065{]} & {[}6.192e-05, 0.0001168{]} & \textless{} 0.0001 \\
A2 - A0 & return\_acf1 & -0.1101 & -3.67 & {[}-0.1206, -0.09987{]} &
{[}-0.1266, -0.09427{]} & \textless{} 0.0001 \\
A2 - A0 & max\_drawdown & -0.004808 & -1.86 & {[}-0.005789, -0.003958{]}
& {[}-0.006307, -0.00357{]} & \textless{} 0.0001 \\
A3 - A0 & vci & 0.02751 & 1.63 & {[}0.02169, 0.03358{]} & {[}0.01825,
0.03658{]} & \textless{} 0.0001 \\
A3 - A0 & realised\_volatility & -0.000151 & -2.67 & {[}-0.0001705,
-0.0001311{]} & {[}-0.0001814, -0.0001195{]} & \textless{} 0.0001 \\
A3 - A0 & return\_acf1 & 0.01642 & 0.603 & {[}0.006676, 0.02584{]} &
{[}0.001701, 0.03075{]} & \textless{} 0.0001 \\
A3 - A0 & max\_drawdown & 0.002757 & 0.732 & {[}0.001441, 0.004085{]} &
{[}0.0008297, 0.004832{]} & \textless{} 0.0001 \\
A1 - A2 & vci & -0.02568 & -0.373 & {[}-0.04646, 0.001704{]} &
{[}-0.05028, 0.01248{]} & 0.377 \\
A1 - A2 & realised\_volatility & -0.000867 & -4.66 & {[}-0.000929,
-0.0007987{]} & {[}-0.0009543, -0.0007547{]} & \textless{} 0.0001 \\
A1 - A2 & return\_acf1 & 0.4829 & 4.81 & {[}0.4528, 0.5228{]} &
{[}0.4455, 0.5446{]} & \textless{} 0.0001 \\
A1 - A2 & max\_drawdown & 0.1947 & 3.32 & {[}0.1754, 0.2169{]} &
{[}0.1687, 0.2297{]} & \textless{} 0.0001 \\
\end{longtable}

\begin{figure}
\centering
\pandocbounded{\includesvg[keepaspectratio]{figures/EXP-001_vci_paired_diff.svg}}
\caption{Figure 1. EXP-001 paired VCI differences against A0, one dot
per seed (30 seeds per contrast); red bar = mean difference. VCI is
dimensionless; dt = 50 steps.}
\end{figure}

\begin{figure}
\centering
\pandocbounded{\includesvg[keepaspectratio]{figures/EXP-001_mid_paths.svg}}
\caption{Figure 2. EXP-001 mid-price path of seed 1001 in each arm
(price units against event time).}
\end{figure}

\subsection{5.2 EXP-002: liquidity supply and
spread}\label{exp-002-liquidity-supply-and-spread}

In the ZI-only market, the mean quoted spread at low liquidity (ZI
market-order probability 0.35) exceeds that at high liquidity (0.1) by
1.632 ticks, 95\% interval {[}1.624, 1.641{]}, \texttt{d\_z} = 69.3, 30
pairs: H1 is supported. The spread falls monotonically (2.923, 1.792 and
1.291 ticks at low, medium and high liquidity) and depth rises.

The pre-declared degenerate arm low-ZI+TF was flagged in 30 of 30 runs
(one-sided book in more than 5\% of steps; mean maximum drawdown 0.9996,
the mid falling to the one-tick floor). Under the protocol it is a model
failure for that arm; it was reported in full and excluded from every
contrast, as fixed before the freeze (amendment A1).

The secondary interaction
\texttt{{[}VCI(med-ZI+TF)\ -\ VCI(med-ZI){]}\ -\ {[}VCI(high-ZI+TF)\ -\ VCI(high-ZI){]}}
is 0.0555, Holm interval {[}0.03614, 0.08209{]}. Because the
medium-liquidity TF effect itself has an interval including 0 (it is the
EXP-001 contrast), this means the TF effect on VCI is lower at high
liquidity, where high-ZI+TF has a mean VCI of -0.03352, not that TF
raise clustering at medium liquidity. A negative VCI is itself unusual:
in high-ZI+TF the mean spread is 1.132 ticks, so the mid moves in
half-tick steps and alternation can dominate the absolute returns. The
interaction may therefore reflect discreteness at a spread close to one
tick rather than an interplay between liquidity and trend feedback
(review pass R2-6).

\textbf{Table 4.} EXP-002 arm summary, mean (SD) over 30 runs. Source:
\texttt{results/confirmatory/EXP-002/tables/arm\_summary.csv}.

\begin{longtable}[]{@{}
  >{\raggedright\arraybackslash}p{(\linewidth - 12\tabcolsep) * \real{0.1429}}
  >{\raggedright\arraybackslash}p{(\linewidth - 12\tabcolsep) * \real{0.1429}}
  >{\raggedright\arraybackslash}p{(\linewidth - 12\tabcolsep) * \real{0.1429}}
  >{\raggedright\arraybackslash}p{(\linewidth - 12\tabcolsep) * \real{0.1429}}
  >{\raggedright\arraybackslash}p{(\linewidth - 12\tabcolsep) * \real{0.1429}}
  >{\raggedright\arraybackslash}p{(\linewidth - 12\tabcolsep) * \real{0.1429}}
  >{\raggedright\arraybackslash}p{(\linewidth - 12\tabcolsep) * \real{0.1429}}@{}}
\toprule\noalign{}
\begin{minipage}[b]{\linewidth}\raggedright
Arm
\end{minipage} & \begin{minipage}[b]{\linewidth}\raggedright
Runs
\end{minipage} & \begin{minipage}[b]{\linewidth}\raggedright
Flagged
\end{minipage} & \begin{minipage}[b]{\linewidth}\raggedright
Mean spread (ticks)
\end{minipage} & \begin{minipage}[b]{\linewidth}\raggedright
Depth at best (units)
\end{minipage} & \begin{minipage}[b]{\linewidth}\raggedright
VCI
\end{minipage} & \begin{minipage}[b]{\linewidth}\raggedright
Max drawdown
\end{minipage} \\
\midrule\noalign{}
\endhead
\bottomrule\noalign{}
\endlastfoot
high-ZI & 30 & 0 & 1.291 (0.01143) & 120.8 (6.762) & 0.0298 (0.01247) &
0.004952 (0.001149) \\
high-ZI+TF & 30 & 0 & 1.132 (0.005144) & 149.8 (1.829) & -0.03352
(0.004799) & 0.1323 (0.01738) \\
low-ZI & 30 & 0 & 2.923 (0.02446) & 41.59 (1.284) & 0.02165 (0.009086) &
0.02138 (0.005369) \\
low-ZI+TF & 30 & 30 & 3.806 (0.04857) & 14.68 (0.8286) & 0.3689
(0.08861) & 0.9996 (9.863e-06) \\
med-ZI & 30 & 0 & 1.792 (0.01671) & 71.1 (1.992) & 0.02564 (0.008298) &
0.01051 (0.002913) \\
med-ZI+TF & 30 & 0 & 1.3 (0.02458) & 109.4 (5.013) & 0.01781 (0.06483) &
0.2004 (0.05875) \\
\end{longtable}

\subsection{5.3 EXP-003: response to an exogenous sell
shock}\label{exp-003-response-to-an-exogenous-sell-shock}

The post-shock maximum drawdown over 5,000 steps (fraction of the
pre-shock mid) is larger with trend followers: S-ZI+TF minus S-ZI =
0.02545, 95\% interval {[}0.01976, 0.03148{]}, \texttt{d\_z} = 1.52; H1
is supported. The secondary mean-reverter contrast is -0.0007314, Holm
interval {[}-0.00159, 0.0001319{]} (Holm p 0.0932). In 7 of 30 TF runs
the mid did not return within 1\% of the pre-shock level within 20,000
steps (censored); in the ZI-only and MR markets the mid stayed within
the band in at least half of the runs (Kaplan-Meier median 0). Spread
widening, a declared secondary outcome, was not measured, because the
protocol gives no formula and no metric is defined after the data (owner
decision, D-015).

\textbf{Table 5.} EXP-003 recovery (descriptive) and bid depth at the
shock. Source: \texttt{results/confirmatory/EXP-003/analysis.json}.

\begin{longtable}[]{@{}
  >{\raggedright\arraybackslash}p{(\linewidth - 10\tabcolsep) * \real{0.1667}}
  >{\raggedright\arraybackslash}p{(\linewidth - 10\tabcolsep) * \real{0.1667}}
  >{\raggedright\arraybackslash}p{(\linewidth - 10\tabcolsep) * \real{0.1667}}
  >{\raggedright\arraybackslash}p{(\linewidth - 10\tabcolsep) * \real{0.1667}}
  >{\raggedright\arraybackslash}p{(\linewidth - 10\tabcolsep) * \real{0.1667}}
  >{\raggedright\arraybackslash}p{(\linewidth - 10\tabcolsep) * \real{0.1667}}@{}}
\toprule\noalign{}
\begin{minipage}[b]{\linewidth}\raggedright
Arm
\end{minipage} & \begin{minipage}[b]{\linewidth}\raggedright
Runs
\end{minipage} & \begin{minipage}[b]{\linewidth}\raggedright
Recovered
\end{minipage} & \begin{minipage}[b]{\linewidth}\raggedright
Censored at 20,000 steps
\end{minipage} & \begin{minipage}[b]{\linewidth}\raggedright
Kaplan-Meier median (steps)
\end{minipage} & \begin{minipage}[b]{\linewidth}\raggedright
Bid depth at shock, mean (SD)
\end{minipage} \\
\midrule\noalign{}
\endhead
\bottomrule\noalign{}
\endlastfoot
S-ZI & 30 & 30 & 0 & 0 & 1551 (94.86) \\
S-ZI+TF & 30 & 23 & 7 & 1,537 & 1285 (564.9) \\
S-ZI+MR & 30 & 30 & 0 & 0 & 1191 (123.2) \\
\end{longtable}

\subsection{5.4 EXP-004: sensitivity
analyses}\label{exp-004-sensitivity-analyses}

\textbf{Table 6.} EXP-004 components (30 pairs each). Source:
\texttt{results/confirmatory/EXP-004/analysis.json}.

\begin{longtable}[]{@{}
  >{\raggedright\arraybackslash}p{(\linewidth - 12\tabcolsep) * \real{0.1429}}
  >{\raggedright\arraybackslash}p{(\linewidth - 12\tabcolsep) * \real{0.1429}}
  >{\raggedright\arraybackslash}p{(\linewidth - 12\tabcolsep) * \real{0.1429}}
  >{\raggedright\arraybackslash}p{(\linewidth - 12\tabcolsep) * \real{0.1429}}
  >{\raggedright\arraybackslash}p{(\linewidth - 12\tabcolsep) * \real{0.1429}}
  >{\raggedright\arraybackslash}p{(\linewidth - 12\tabcolsep) * \real{0.1429}}
  >{\raggedright\arraybackslash}p{(\linewidth - 12\tabcolsep) * \real{0.1429}}@{}}
\toprule\noalign{}
\begin{minipage}[b]{\linewidth}\raggedright
Component
\end{minipage} & \begin{minipage}[b]{\linewidth}\raggedright
Contrast
\end{minipage} & \begin{minipage}[b]{\linewidth}\raggedright
Metric
\end{minipage} & \begin{minipage}[b]{\linewidth}\raggedright
Pairs
\end{minipage} & \begin{minipage}[b]{\linewidth}\raggedright
Mean diff.
\end{minipage} & \begin{minipage}[b]{\linewidth}\raggedright
d\_z
\end{minipage} & \begin{minipage}[b]{\linewidth}\raggedright
95\% interval
\end{minipage} \\
\midrule\noalign{}
\endhead
\bottomrule\noalign{}
\endlastfoot
placebo: EXP-001 A1 minus EXP-004 A4 & A1 - A4 & vci & 30 & -0.02534 &
-0.402 & {[}-0.04414, -0.0002837{]} \\
descriptive: A4 minus EXP-001 A0 & A4 - A0 & vci & 30 & 0.01751 & 1.54 &
{[}0.0135, 0.02148{]} \\
tick 2: A1-tick2 minus A0-tick2 & A1-tick2 - A0-tick2 & vci & 30 &
-0.005427 & -0.0656 & {[}-0.03144, 0.02572{]} \\
sampling interval vci\_dt10 & A1 - A0 & vci\_dt10 & 30 & 0.02481 & 0.728
& {[}0.0153, 0.03842{]} \\
sampling interval vci\_dt200 & A1 - A0 & vci\_dt200 & 30 & 0.0915 &
0.608 & {[}0.04287, 0.1489{]} \\
exclusion of flagged runs and partners & A1 - A0 & vci & 30 & -0.007825
& -0.121 & {[}-0.0268, 0.01785{]} \\
\end{longtable}

\emph{Placebo.} A4 minus A0 is 0.01751 {[}0.0135, 0.02148{]}:
price-independent aggressive flow alone raises VCI at \texttt{dt} = 50.
A1 minus A4 is -0.02534 {[}-0.04414, -0.0002837{]}: at this interval
there is no evidence that the trend rule adds clustering beyond
aggressiveness-matched coin-flip flow. The point estimate is negative,
but the interval's upper end is close to 0 and the six EXP-004
components carry no multiplicity adjustment, so ``less clustering than
the placebo'' is not a robust statement (review pass R2-3). Because
EXP-001 did not support H1, the preregistered attribution rule is not
triggered.

\emph{Sampling interval.} The A1 - A0 contrast is 0.02481 {[}0.0153,
0.03842{]} at \texttt{dt} = 10 and 0.0915 {[}0.04287, 0.1489{]} at
\texttt{dt} = 200. \emph{Tick size.} At tick 2 it is -0.005427
{[}-0.03144, 0.02572{]}.

\emph{Frozen robustness rule.} The rule calls EXP-001 robust only if
every \texttt{dt} and tick component has the same sign as EXP-001 and an
interval excluding 0, and invalidated by a sign reversal. With ``same
sign'' read as the sign of the EXP-001 point estimate (D-013), the
\texttt{dt} components are reversals with intervals excluding 0, and the
rule returns \textbf{invalidated (a sign reversal with an interval
excluding 0)}. The owner decided that this is reported exactly as the
rule gives it, with the note that the primary estimate is near zero, so
its sign is weakly determined (D-015).

\begin{figure}
\centering
\pandocbounded{\includesvg[keepaspectratio]{figures/EXP-004_frozen_contrasts.svg}}
\caption{Figure 3. EXP-004 components: mean paired VCI difference over
30 seeds (red dot) and 95\% percentile bootstrap interval (blue line);
dashed line at 0.}
\end{figure}

\subsection{5.5 EXP-005: simulated and real-market statistics
(FI-2010)}\label{exp-005-simulated-and-real-market-statistics-fi-2010}

\emph{(a) Held-out check.} On the held-out day 10, VCI at the matched
interval (\texttt{h} = 5 columns, 50 events) is -0.001241, 0.03691,
0.05438, 0.03336 and 0.07877 in the five stocks. The frozen hypothesis
(VCI positive in every stock) is not supported, because one stock is
just below 0. Single stock-days (1,888 to 13,378 columns) give noisy VCI
estimates.

\emph{(b) Placement.} Against the range of 37 reference blocks of days 1
to 9 (Table 7), every arm's VCI interval lies inside or overlaps the
real range, which is wide (0.0004353 to 0.168). Realised volatility at
the primary interval is above the real range for A0, A2 and A3 and
inside it for A1; the relative spread is inside for A0, A1 and A3 and
above for A2. This placement is descriptive, depends on a crude clock
match (one agent action against one tenth of an FI-2010 row) and on the
tick-to-price ratio, and is never a validation of the model.

\textbf{Table 7.} EXP-005 placement at the primary sampling. Source:
\texttt{results/confirmatory/EXP-005/comparison.json}.

\begin{longtable}[]{@{}
  >{\raggedright\arraybackslash}p{(\linewidth - 12\tabcolsep) * \real{0.1429}}
  >{\raggedright\arraybackslash}p{(\linewidth - 12\tabcolsep) * \real{0.1429}}
  >{\raggedright\arraybackslash}p{(\linewidth - 12\tabcolsep) * \real{0.1429}}
  >{\raggedright\arraybackslash}p{(\linewidth - 12\tabcolsep) * \real{0.1429}}
  >{\raggedright\arraybackslash}p{(\linewidth - 12\tabcolsep) * \real{0.1429}}
  >{\raggedright\arraybackslash}p{(\linewidth - 12\tabcolsep) * \real{0.1429}}
  >{\raggedright\arraybackslash}p{(\linewidth - 12\tabcolsep) * \real{0.1429}}@{}}
\toprule\noalign{}
\begin{minipage}[b]{\linewidth}\raggedright
Arm
\end{minipage} & \begin{minipage}[b]{\linewidth}\raggedright
Statistic
\end{minipage} & \begin{minipage}[b]{\linewidth}\raggedright
Arm mean
\end{minipage} & \begin{minipage}[b]{\linewidth}\raggedright
95\% interval
\end{minipage} & \begin{minipage}[b]{\linewidth}\raggedright
FI-2010 median
\end{minipage} & \begin{minipage}[b]{\linewidth}\raggedright
FI-2010 {[}min, max{]}
\end{minipage} & \begin{minipage}[b]{\linewidth}\raggedright
Placement
\end{minipage} \\
\midrule\noalign{}
\endhead
\bottomrule\noalign{}
\endlastfoot
A0 & vci & 0.02564 & {[}0.02274, 0.02855{]} & 0.0481 & {[}0.0004353,
0.168{]} & inside \\
A0 & realised\_volatility & 0.001302 & {[}0.001289, 0.001314{]} &
0.0004999 & {[}0.0002354, 0.001041{]} & above \\
A0 & mean\_relative\_spread & 0.001791 & {[}0.001785, 0.001797{]} &
0.001675 & {[}0.001098, 0.00234{]} & inside \\
A1 & vci & 0.01781 & {[}-0.001045, 0.04381{]} & 0.0481 & {[}0.0004353,
0.168{]} & overlapping \\
A1 & realised\_volatility & 0.0005241 & {[}0.000463, 0.0005953{]} &
0.0004999 & {[}0.0002354, 0.001041{]} & inside \\
A1 & mean\_relative\_spread & 0.001394 & {[}0.001244, 0.001562{]} &
0.001675 & {[}0.001098, 0.00234{]} & inside \\
A2 & vci & 0.0435 & {[}0.03731, 0.04942{]} & 0.0481 & {[}0.0004353,
0.168{]} & inside \\
A2 & realised\_volatility & 0.001391 & {[}0.001378, 0.001405{]} &
0.0004999 & {[}0.0002354, 0.001041{]} & above \\
A2 & mean\_relative\_spread & 0.002765 & {[}0.002755, 0.002776{]} &
0.001675 & {[}0.001098, 0.00234{]} & above \\
A3 & vci & 0.05315 & {[}0.04789, 0.05885{]} & 0.0481 & {[}0.0004353,
0.168{]} & inside \\
A3 & realised\_volatility & 0.001151 & {[}0.001137, 0.001164{]} &
0.0004999 & {[}0.0002354, 0.001041{]} & above \\
A3 & mean\_relative\_spread & 0.001955 & {[}0.001944, 0.001966{]} &
0.001675 & {[}0.001098, 0.00234{]} & inside \\
\end{longtable}

\subsection{5.6 Exploratory M4 checks (not
confirmatory)}\label{exploratory-m4-checks-not-confirmatory}

These checks were run after the confirmatory analysis, are labelled
exploratory, and do not change any verdict (journal J-0011, D-016,
\texttt{results/exploratory/M4-SENS/}). New simulations used development
seeds 101 to 130 only, with the frozen parameter values; 180 runs, 58.4
seconds of wall time in the final execution, two workers. The script was
executed three times; each execution overwrote the previous output with
identical values for every quantity the earlier versions computed, and
added the E6 and signed-autocorrelation columns (J-0011).

\emph{E1. Sampling-interval grid.} On the development seeds, the
confirmatory pattern recurs: A1 - A0 is -0.006956 {[}-0.02487,
0.01528{]} at \texttt{dt} = 50 and positive at 10 and 200. Across
\texttt{dt} from 5 to 500 steps (Table 8, Figure 4), the placebo
contrast A4 - A0 is positive with intervals excluding 0 up to
\texttt{dt} = 100 and near zero from 200 on. A1 - A0 is positive at 5
and 10 steps, has intervals including 0 from 20 to 100, and is large at
200 and 500 (0.2334 {[}0.1693, 0.2994{]} at 500). The contrast A1 - A4
is negative at 10 to 50 steps and positive at 200 and 500 (0.09572
{[}0.04667, 0.1486{]} and 0.2363 {[}0.1737, 0.3021{]}).

The lag-1 autocorrelation of signed returns, recorded at each
\texttt{dt} after the second review pass (R2-1), changes how these
numbers can be read. In A1 it averages -0.013 at \texttt{dt} = 50, 0.653
at 200 and 0.8921 at 500: at long intervals the trend-follower market
moves in runs, and a run of same-signed returns makes absolute returns
autocorrelated whether or not volatility clusters. In A0 it is negative
at every interval (-0.1691 at \texttt{dt} = 5 to -0.4315 at 200), the
signature of a mid bouncing inside the spread, and in A4 it is similar
(-0.3316 at 50). The defensible reading is therefore narrow: at short
intervals, adding aggressive order flow of either kind changes VCI in a
market dominated by bounce; at long intervals, trend followers produce
trending runs that VCI cannot distinguish from clustering. Whether trend
feedback produces volatility clustering in the stylized-fact sense is
not answered by VCI in this model. These checks were designed after the
confirmatory results were seen and are hypothesis-generating only
(R2-7).

\textbf{Table 8.} EXPLORATORY: VCI contrasts by sampling interval,
development seeds 101 to 130 (30 pairs each). Source:
\texttt{results/exploratory/M4-SENS/m4\_exploratory.json}.

\begin{longtable}[]{@{}
  >{\raggedright\arraybackslash}p{(\linewidth - 8\tabcolsep) * \real{0.2000}}
  >{\raggedright\arraybackslash}p{(\linewidth - 8\tabcolsep) * \real{0.2000}}
  >{\raggedright\arraybackslash}p{(\linewidth - 8\tabcolsep) * \real{0.2000}}
  >{\raggedright\arraybackslash}p{(\linewidth - 8\tabcolsep) * \real{0.2000}}
  >{\raggedright\arraybackslash}p{(\linewidth - 8\tabcolsep) * \real{0.2000}}@{}}
\toprule\noalign{}
\begin{minipage}[b]{\linewidth}\raggedright
dt (steps)
\end{minipage} & \begin{minipage}[b]{\linewidth}\raggedright
A1 - A0 (trend)
\end{minipage} & \begin{minipage}[b]{\linewidth}\raggedright
A4 - A0 (placebo)
\end{minipage} & \begin{minipage}[b]{\linewidth}\raggedright
A1 - A4 (feedback beyond aggressiveness)
\end{minipage} & \begin{minipage}[b]{\linewidth}\raggedright
A2 - A0 (mean reversion)
\end{minipage} \\
\midrule\noalign{}
\endhead
\bottomrule\noalign{}
\endlastfoot
5 & 0.03437 {[}0.02843, 0.04242{]} & 0.03667 {[}0.03395, 0.03941{]} &
-0.002297 {[}-0.009569, 0.006964{]} & -0.01909 {[}-0.02092,
-0.01726{]} \\
10 & 0.02067 {[}0.01207, 0.03211{]} & 0.041 {[}0.03797, 0.0439{]} &
-0.02033 {[}-0.02961, -0.008532{]} & -0.01052 {[}-0.01384,
-0.007346{]} \\
20 & 0.001766 {[}-0.01146, 0.01895{]} & 0.03696 {[}0.0333, 0.04061{]} &
-0.0352 {[}-0.0477, -0.01942{]} & 0.007182 {[}0.003053, 0.01134{]} \\
50 & -0.006956 {[}-0.02487, 0.01528{]} & 0.0213 {[}0.0158, 0.02678{]} &
-0.02826 {[}-0.0466, -0.005951{]} & 0.0177 {[}0.01102, 0.02442{]} \\
100 & 0.02324 {[}-0.01016, 0.0626{]} & 0.009015 {[}0.001338, 0.01607{]}
& 0.01423 {[}-0.0175, 0.05241{]} & 0.02494 {[}0.01771, 0.03232{]} \\
200 & 0.09813 {[}0.04939, 0.1508{]} & 0.002401 {[}-0.007938, 0.01207{]}
& 0.09572 {[}0.04667, 0.1486{]} & 0.0218 {[}0.01033, 0.03407{]} \\
500 & 0.2334 {[}0.1693, 0.2994{]} & -0.002919 {[}-0.01283, 0.006592{]} &
0.2363 {[}0.1737, 0.3021{]} & 0.01441 {[}0.003402, 0.02573{]} \\
\end{longtable}

\begin{figure}
\centering
\pandocbounded{\includesvg[keepaspectratio]{figures/M4_exploratory_dt_grid.svg}}
\caption{Figure 4. EXPLORATORY: A1 minus A0 and A1 minus A4 VCI
contrasts across sampling intervals of 5 to 500 steps, 30 development
seeds, 95\% percentile bootstrap intervals.}
\end{figure}

\emph{E2. Unshocked control window for EXP-003.} The EXP-003 design has
no unshocked TF market, so part of the post-shock drawdown could be an
ordinary trend excursion. On the development seeds, the drawdown formula
applied at step 50,000 in unshocked runs gives A1 - A0 = 0.007912
{[}0.004752, 0.01107{]}; the shocked contrast replicates (0.02408
{[}0.01689, 0.03223{]}); the difference-in-differences
\texttt{{[}S-ZI+TF\ -\ A1{]}\ -\ {[}S-ZI\ -\ A0{]}} is 0.01616
{[}0.009609, 0.02483{]} (Table 9). On these seeds about one third of the
shocked contrast is matched by ordinary TF excursions, and the remainder
is a larger shock response in the TF market. The decomposition is
approximate (a drawdown is a minimum over the window, so effects need
not add), and it uses different seeds from EXP-003.

\textbf{Table 9.} EXPLORATORY: post-shock window drawdown with and
without the shock, development seeds 101 to 130.

\begin{longtable}[]{@{}
  >{\raggedright\arraybackslash}p{(\linewidth - 8\tabcolsep) * \real{0.2000}}
  >{\raggedright\arraybackslash}p{(\linewidth - 8\tabcolsep) * \real{0.2000}}
  >{\raggedright\arraybackslash}p{(\linewidth - 8\tabcolsep) * \real{0.2000}}
  >{\raggedright\arraybackslash}p{(\linewidth - 8\tabcolsep) * \real{0.2000}}
  >{\raggedright\arraybackslash}p{(\linewidth - 8\tabcolsep) * \real{0.2000}}@{}}
\toprule\noalign{}
\begin{minipage}[b]{\linewidth}\raggedright
Contrast (window drawdown from step 50,000, fraction of the step-49,999
mid)
\end{minipage} & \begin{minipage}[b]{\linewidth}\raggedright
Pairs
\end{minipage} & \begin{minipage}[b]{\linewidth}\raggedright
Mean diff.
\end{minipage} & \begin{minipage}[b]{\linewidth}\raggedright
d\_z
\end{minipage} & \begin{minipage}[b]{\linewidth}\raggedright
95\% interval
\end{minipage} \\
\midrule\noalign{}
\endhead
\bottomrule\noalign{}
\endlastfoot
shocked: S-ZI+TF - S-ZI (replicates the EXP-003 primary on development
seeds) & 30 & 0.02408 & 1.1 & {[}0.01689, 0.03223{]} \\
unshocked window: A1 - A0 (ordinary TF excursion, no shock) & 30 &
0.007912 & 0.875 & {[}0.004752, 0.01107{]} \\
shock response, TF market: S-ZI+TF - A1 & 30 & 0.01913 & 0.88 &
{[}0.01263, 0.02781{]} \\
shock response, ZI market: S-ZI - A0 & 30 & 0.002969 & 2.7 &
{[}0.002569, 0.003335{]} \\
difference-in-differences: {[}S-ZI+TF - A1{]} - {[}S-ZI - A0{]} & 30 &
0.01616 & 0.74 & {[}0.009609, 0.02483{]} \\
\end{longtable}

\emph{E3 to E6} (committed confirmatory outputs only). Giving every
contrast tied at bootstrap p = 0 the most conservative Holm level leaves
every interval excluding 0 (supplement S4). Reading the EXP-004 rule
with the sign of H1 as reference turns ``invalidated'' into
``qualified'' (supplement S4). Across A1 seeds, VCI is nearly unrelated
to the run's maximum drawdown (Spearman 0.04605). In EXP-003 the shock
was always fully filled and no run had a one-sided book; within the TF
arm the drawdown is nearly unrelated to bid depth at the shock (Spearman
0.07141).

\section{6. Discussion}\label{discussion}

\subsection{6.1 What the confirmatory evidence
supports}\label{what-the-confirmatory-evidence-supports}

Within this model and parameter set, replacing 30\% of random-order
agents by trend followers did not detectably change volatility
clustering at the preregistered sampling interval of 50 steps. The sign
of the effect depends on the interval: positive at 10 and 200 steps,
near zero at 50. The frozen robustness rule therefore declares the
primary interpretation invalidated, and the only defensible confirmatory
statement is the negative one: \textbf{this study does not show that
trend followers raise volatility clustering at the preregistered
horizon, and the answer depends on the horizon.} The interval dependence
was visible before the freeze: the exploratory pilot on development
seeds 1 to 20 gave A1 - A0 = -0.003781 {[}-0.02328, 0.02025{]} at
\texttt{dt} = 50 and positive contrasts at 10 and 200 (0.01781 and
0.1319; J-0003). It was reported to the owner as a reason to keep
EXP-004 mandatory (D-007 item 6), and the primary interval of 50 steps
was kept unchanged. The confirmatory runs reproduced the pilot pattern.
The exploratory checks add a reason to distrust any single-interval VCI
comparison in this model: at short intervals VCI moves with bid-ask
bounce, and at long intervals it moves with the trend runs that trend
followers create.

Two secondary findings are supported with large standardised effects.
More random-order limit supply narrows the spread, which is close to
mechanical here (a lower market-order probability means more limit
orders per step) and agrees with the random-order theory
{[}@Smith2003{]}. Trend followers deepen the drawdown after a large
sell; the exploratory difference-in-differences suggests that ordinary
excursions account for part, not all, of this.

\subsection{6.2 Competing explanations}\label{competing-explanations}

The EXP-001 protocol names three competing explanations; Table 10 states
which results address them.

\begin{longtable}[]{@{}
  >{\raggedright\arraybackslash}p{(\linewidth - 4\tabcolsep) * \real{0.3333}}
  >{\raggedright\arraybackslash}p{(\linewidth - 4\tabcolsep) * \real{0.3333}}
  >{\raggedright\arraybackslash}p{(\linewidth - 4\tabcolsep) * \real{0.3333}}@{}}
\toprule\noalign{}
\begin{minipage}[b]{\linewidth}\raggedright
Competing explanation (EXP-001 protocol section 3)
\end{minipage} & \begin{minipage}[b]{\linewidth}\raggedright
Addressed by
\end{minipage} & \begin{minipage}[b]{\linewidth}\raggedright
What the results say
\end{minipage} \\
\midrule\noalign{}
\endhead
\bottomrule\noalign{}
\endlastfoot
1. Aggressiveness: TF send market orders at a different rate & EXP-004
placebo (confirmatory); E1 (exploratory) & Aggressive flow alone raises
VCI at \texttt{dt} = 50 (A4 - A0 0.01751) and there is no evidence that
the trend rule adds to it at that interval (A1 - A4 -0.02534, upper end
close to 0). In E1, the placebo effect is confined to short
intervals. \\
2. Liquidity removal: each TF replaces a ZI liquidity provider & A4 has
the same ZI count as A1 (confirmatory); EXP-002 interaction (secondary)
& A1 - A4 holds liquidity providers fixed, so the A1 - A4 contrast is
free of this confound; the interaction shows the TF effect on VCI
changes with liquidity. Not addressed: the depth difference between A1
and A0 itself. \\
3. Sampling artefact of \texttt{dt} & EXP-004 \texttt{dt} components
(confirmatory); E1 (exploratory) & The effect is horizon-dependent. This
is the reason the rule returns ``invalidated''. \\
Not named in the protocol: trend runs inflate the autocorrelation of
absolute returns & E1 signed-return autocorrelation; E5 (exploratory) &
Holds at long intervals: A1's signed-return autocorrelation is 0.653 at
\texttt{dt} = 200 and 0.8921 at 500, so the long-horizon VCI increase
reflects runs at least as much as clustering. At \texttt{dt} = 50 it is
near zero, and A1 VCI is nearly unrelated to the run's maximum drawdown
(Spearman 0.04605). \\
Not named in the protocol: bid-ask bounce shapes VCI at short intervals
& Committed arm summaries; E1 & The baseline's signed-return
autocorrelation is strongly negative (A0 -0.3786 at \texttt{dt} = 50 in
the confirmatory runs; A1 -0.00584). A0 - A1 differences in VCI mix a
change in bounce with any change in clustering. Not resolved by this
design. \\
Not named in the protocol: the EXP-002 interaction reflects discreteness
& Committed arm summaries & Cannot be excluded: high-ZI+TF has a
negative VCI (-0.03352) at a spread of 1.132 ticks. \\
EXP-003: bid depth at the shock differs by arm & E6 (exploratory,
committed runs) & Within the TF arm drawdown is nearly unrelated to
depth at the shock (Spearman 0.07141); the arm means differ (1,285
against 1,551 units). \\
EXP-003: the shock consumes one side completely & E6 & The shock was
fully filled in 90 of 90 runs and no run had a one-sided book. \\
EXP-003: ordinary TF excursions overlap the window & E2 (exploratory,
development seeds) & About one third of the effect on those seeds; the
difference-in-differences remains positive with an interval excluding
0. \\
EXP-002: order-lifetime effects and one-sided episodes & Not addressed
(the halved-lifetime sensitivity named in the protocol was not run) &
The spread result remains largely mechanical. \\
\end{longtable}

\subsection{6.3 Relation to prior work}\label{relation-to-prior-work}

The result does not reproduce, at the primary interval, the attribution
of clustering to chartists reported by @Chiarella2009, but the two
models differ (their agents include fundamentalists, their trading and
sampling differ, and only their abstract was read), so this is not a
contradiction of that work. The placebo result and the exploratory sweep
give the more general lesson: in an order-driven market, changing the
share of aggressive orders changes the autocorrelation of absolute
returns, and trend rules create runs that the same statistic reads as
clustering. A chartist-versus-noise comparison at one sampling interval,
without an aggressiveness-matched control and without a check of
signed-return dependence, can therefore confound three mechanisms. The
FI-2010 placement shows only that the simulated statistics are of the
same order of magnitude as one real market under a crude clock match; it
is not a test of the model.

\subsection{6.4 Review of the D-013 analysis
choices}\label{review-of-the-d-013-analysis-choices}

The owner asked for the D-013 choices to be reviewed in M4 (D-015). (1)
\emph{Holm intervals and tie ordering.} Giving every tied contrast the
most conservative level (\texttt{1\ -\ 0.05\ /\ m}) leaves every Holm
interval that excluded 0 still excluding 0 in EXP-001 (24 tied of 30)
and EXP-002 (16 tied of 16); the EXP-002 interaction widens from
{[}0.03614, 0.08209{]} to {[}0.0307, 0.09768{]}. No decision depends on
the ordering, as D-014 stated. (2, 3) \emph{Family definitions.} The
secondary families are reported in full; the alternative EXP-002 reading
(all ten pairs of the five arms) was not computed. (4) \emph{Spread
widening} stays not measured (D-015). (5) \emph{Degenerate-run
sensitivity} is moot: no run in a primary contrast was flagged. (6)
\emph{Sign reference of the robustness rule.} This choice decides the
EXP-004 verdict. Read with the sign of H1 (positive), which is the
direction the rule was written to protect, both \texttt{dt} components
agree in sign and exclude 0 and the tick-2 component includes 0, so the
rule would return ``qualified'' rather than ``invalidated''. The frozen
verdict stands (D-015); both readings agree that the EXP-001 conclusion
depends on the horizon.

\subsection{6.5 Practical significance}\label{practical-significance}

None for real markets. The study's use is methodological: a small,
reproducible demonstration that a widely used comparison (chartists
against noise traders) needs an aggressiveness-matched control, a
horizon sweep and a clustering measure that is not driven by
signed-return dependence before its result can be attributed to
feedback.

\section{7. Limitations and threats to
validity}\label{limitations-and-threats-to-validity}

\begin{itemize}
\tightlist
\item
  \textbf{Construct.} VCI is one statistic (mean ACF of absolute
  returns, lags 1 to 10). Its value changes with the sampling interval,
  and event time is not clock time. In this model it is not a clean
  measure of volatility clustering: at short intervals it is shaped by
  the strongly negative signed-return autocorrelation of a bouncing mid,
  and at long intervals in trend-follower markets by trend runs
  (sections 5.6 and 6.2). Other clustering measures (GARCH fits,
  long-memory exponents, autocorrelation of absolute residuals after
  removing signed dependence) were not used.
\item
  \textbf{Internal.} Replacing ZI by TF also changes depth and spread;
  the placebo matches order types, sizes and activity but not the serial
  correlation of direction, which is the feedback itself. The EXP-002
  liquidity manipulation changes aggressive flow and limit supply
  together. EXP-003 has no unshocked control in the confirmatory design
  (addressed only exploratorily).
\item
  \textbf{Statistical.} A1's VCI varies widely across seeds, so the
  primary test had low power for the effect that occurred (observed
  \texttt{\textbar{}d\_z\textbar{}} about 0.12 against a planned 0.53).
  A null result is not evidence of no effect. Holm intervals among
  contrasts tied at p = 0 depend on input order (D-014); decisions do
  not. The exploratory M4 checks were run after the confirmatory outputs
  were seen, on different seeds, and are hypothesis-generating only.
\item
  \textbf{Analysis choices.} The EXP-004 verdict depends on the sign
  reference chosen in D-013 (section 6.4).
\item
  \textbf{External.} One parameter set, 100 agents, one asset,
  fixed-rule agents that neither learn nor optimise, self-trades
  allowed, a one-tick price floor. The model is not calibrated to any
  market. FI-2010 covers five Helsinki stocks over ten days in 2010; its
  reference blocks may contain overnight returns at unknown positions,
  and the clock and tick matches are crude.
\item
  \textbf{Literature.} Except FI-2010 and @RatliffCrain2023, only
  abstracts or bibliographic records were read; no systematic search for
  placebo controls in agent-based order-book studies was made.
\item
  \textbf{Simulation scope.} Nothing here is a statement about real
  markets, a prediction, trading advice or a claim of profitability.
  Wealth by agent type was recorded descriptively and is not analysed.
\end{itemize}

The full list is in \texttt{research/limitations.md}.

\section{8. Ethics and responsible
research}\label{ethics-and-responsible-research}

The study involves no human participants and no personal data; no ethics
approval was required or is claimed. The only external data, FI-2010, is
public under CC BY 4.0 and is attributed here; it is downloaded by a
script that checks its SHA-256 and is not redistributed in the
repository. To limit misuse, the manuscript makes no statement about
real-market behaviour, trading strategies or profitability, and
simulated or exploratory outcomes are never presented as evidence about
real markets.

\section{9. Conclusion and future
work}\label{conclusion-and-future-work}

Under the frozen protocol, the primary interpretation of this small,
preregistered, paired-seed simulation is invalidated: trend followers
did not detectably change volatility clustering at the preregistered
horizon, and the frozen robustness rule found a sign reversal at other
sampling intervals. As a secondary reading, the same rule taken against
the sign of H1 would return ``qualified''; under either reading the
effect's sign depends on the sampling interval. Price-independent
aggressive order flow raised VCI at that horizon, with no evidence that
trend following added to it. Liquidity supply narrowed the spread, and
trend followers deepened post-shock drawdowns. The feasible next study
is a new preregistered experiment, frozen by the owner, that sweeps the
sampling interval (and the trend lookback) with a placebo arm at each
point, with a clustering measure that is robust to signed-return
dependence, so that the exploratory pattern of section 5.6 can be tested
on fresh seeds. A second extension would vary liquidity without changing
aggressive flow, for example through order lifetime.

\section{10. Data and code
availability}\label{data-and-code-availability}

Repository: \texttt{market-order-book-agents} (public URL:
{[}placeholder, not yet published{]}). Code in \texttt{src/market/};
frozen protocols and configs in \texttt{research/} and
\texttt{experiments/} with their hashes in
\texttt{research/freeze.json}; raw per-run results, analysis outputs and
SHA256SUMS in \texttt{results/confirmatory/}; exploratory outputs in
\texttt{results/exploratory/}; experiment reports in
\texttt{research/experiments/EXP-00N/report.md}; the artifact manifest
and evidence audit in \texttt{research/publication/}. Reproduction
(deterministic given seed and revision):

\begin{verbatim}
/Users/bechir/research/labs/.tools/py/bin/uv sync --frozen --all-groups --all-extras
.venv/bin/python -m pytest
.venv/bin/python -m market.cli run --exp EXP-001 --workers 2     # likewise EXP-002 to EXP-004
.venv/bin/python -m market.frozen_analysis --exp EXP-001         # likewise EXP-002 to EXP-004
.venv/bin/python -m market.cli realdata compare --exp EXP-005 --results results/confirmatory/EXP-001
.venv/bin/python scripts/m4_exploratory.py --workers 2           # exploratory M4 checks
.venv/bin/python scripts/m4_tables.py                            # manuscript tables and figures
.venv/bin/python scripts/m4_audit.py                             # artifact manifest and audit
\end{verbatim}

Real data: FI-2010 {[}@Ntakaris2018{]}, CC BY 4.0, fetched by
\texttt{scripts/download\_fi2010.py} with SHA-256 checks
(\texttt{data/manifest.json}).

\section{11. Funding, competing interests and AI
assistance}\label{funding-competing-interests-and-ai-assistance}

Funding and competing interests: none are recorded in the project files.
The owner confirms both statements before any release.

AI assistance: an AI system (Claude, Anthropic) was used as a research
and development assistant for code, tests, protocol drafting, analysis
scripts, the M4 exploratory checks, the evidence audit and the drafting
of this manuscript, under the owner's direction. The owner approved the
hypotheses, the protocol freeze, the amendments and every decision
recorded in the decision log. Every number in this manuscript is
generated from, or checked by script against, committed results
(\texttt{scripts/m4\_tables.py}, \texttt{scripts/m4\_audit.py}). The
AI's internal review (\texttt{research/publication/review\_pass\_2.md})
is quality control, not peer review. The AI system is not an author.

\section{References}\label{references}

The bibliography is \texttt{references.bib}; every entry was checked
against the page listed in \texttt{research/literature\_review.md}.

\begin{itemize}
\tightlist
\item
  {[}@GodeSunder1993{]} Gode, D. K., and Sunder, S. (1993). Allocative
  efficiency of markets with zero-intelligence traders. \emph{Journal of
  Political Economy} 101(1), 119-137. doi:10.1086/261868. Abstract
  record read.
\item
  {[}@Smith2003{]} Smith, E., Farmer, J. D., Gillemot, L., and
  Krishnamurthy, S. (2003). Statistical theory of the continuous double
  auction. \emph{Quantitative Finance} 3(6), 481-514.
  doi:10.1088/1469-7688/3/6/307. arXiv cond-mat/0210475, abstract read.
\item
  {[}@Farmer2005{]} Farmer, J. D., Patelli, P., and Zovko, I. I. (2005).
  The predictive power of zero intelligence in financial markets.
  \emph{PNAS} 102(11), 2254-2259. doi:10.1073/pnas.0409157102. arXiv
  cond-mat/0309233, abstract read; journal metadata from a search
  result.
\item
  {[}@Chiarella2009{]} Chiarella, C., Iori, G., and Perello, J. (2009).
  The impact of heterogeneous trading rules on the limit order book and
  order flows. \emph{Journal of Economic Dynamics and Control} 33(3),
  525-537. doi:10.1016/j.jedc.2008.08.001. arXiv 0711.3581, abstract
  read.
\item
  {[}@LuxMarchesi1999{]} Lux, T., and Marchesi, M. (1999). Scaling and
  criticality in a stochastic multi-agent model of a financial market.
  \emph{Nature} 397, 498-500. doi:10.1038/17290. Abstract record read.
\item
  {[}@Cont2001{]} Cont, R. (2001). Empirical properties of asset
  returns: stylized facts and statistical issues. \emph{Quantitative
  Finance} 1(2), 223-236. doi:10.1080/713665670. Abstract record read;
  full text not read.
\item
  {[}@RatliffCrain2023{]} Ratliff-Crain, E., Van Oort, C. M., Bagrow,
  J., Koehler, M. T. K., and Tivnan, B. F. (2023). Revisiting Cont's
  stylized facts for modern stock markets. arXiv 2311.07738. Full text
  (HTML, v2) read 2026-11-25.
\item
  {[}@Chakraborti2011{]} Chakraborti, A., Muni Toke, I., Patriarca, M.,
  and Abergel, F. (2011). Econophysics: empirical facts and agent-based
  models. \emph{Quantitative Finance} 11, 991-1012 and 1013-1041. arXiv
  0909.1974, abstract read.
\item
  {[}@Byrd2019{]} Byrd, D., Hybinette, M., and Hybinette Balch, T.
  (2019). ABIDES: towards high-fidelity market simulation for AI
  research. arXiv 1904.12066. Abstract read.
\item
  {[}@Ntakaris2018{]} Ntakaris, A., Magris, M., Kanniainen, J., Gabbouj,
  M., and Iosifidis, A. (2018). Benchmark dataset for mid-price
  forecasting of limit order book data with machine learning methods.
  \emph{Journal of Forecasting} 37(8), 852-866. doi:10.1002/for.2543.
  arXiv 1705.03233 v5, full PDF opened 2026-11-24. Dataset CC BY 4.0.
\end{itemize}

\section{Supplementary appendix}\label{supplementary-appendix}

The supplement (\texttt{supplementary.md}) holds the complete
experimental matrix, the configurations, extended tables, the
exploratory robustness results E3 and E4, links to the five experiment
reports and the list of deviations and owner decisions.

\end{document}
